%% kapa: exam=1 round=25 year=2014 period=1 form=A subject=civil_law subject_ko=민법 seq=1 number=1
%% source: Q-Net 감정평가사 기출문제 (cst00309) artlSeq=5101207 file=extracted/1008/5101207/2014년 제25회 감정평가사 1차시험 기출문제/1교시(A형).hwp
%% conversion: kordoc 4.24.51 -> tools/gen1.py (auto); review=proofread-sonnet; flags=none
\begin{question}{1}
법인의 기관에 관한 설명으로 옳지 않은 것은? (다툼이 있으면 판례에 의함)
\begin{choices}
  \item 정관에 다른 규정이 없는 한, 사원총회의 결의사항은 총회를 소집할 때 미리 적법하게 통지한 사항에 한한다.
  \item 이사 사임의 의사표시가 법인의 대표자에게 도달한 후에는, 특별한 사정이 없으면, 그 의사표시를 철회할 수 없다.
  \item 사원총회 소집권자가 총회 구성원들에게 총회 개최의 연기를 통지할 때에는 총회의 소집과 동일한 방식에 의할 필요가 없다.
  \item 정관이나 사원총회의 결의로 금지하지 않는 한, 이사는 타인으로 하여금 법인의 제반사무 처리를 포괄적으로 대리하게 할 수 있다.
  \item 사원총회의 결의는 법률 또는 정관에 다른 규정이 없으면 사원 과반수의 출석과 출석사원의 결의권의 과반수로써 한다.
\end{choices}
\end{question}
\answer{4}
